\addcontentsline{toc}{chapter}{Bibliography}
\begin{thebibliography}{00}

% Entries are numbered in the order they are first cited in the report (IEEE style).
% Keep this order when adding a reference: put the new \bibitem where it is first cited.

%====================== Chapter 1 : Introduction ======================

\bibitem{who2024neuro}
World Health Organization, ``Over 1 in 3 people affected by neurological conditions, the leading cause of illness and disability worldwide,'' \textit{WHO}, Mar.\ 14, 2024. [Online]. Available: \url{https://www.who.int/news/item/14-03-2024-over-1-in-3-people-affected-by-neurological-conditions--the-leading-cause-of-illness-and-disability-worldwide}. [Accessed: April 10, 2026].

%====================== Chapter 2 : Literature Survey ======================

\bibitem{mullen2015wearable}
T. R. Mullen, C. A. E. Kothe, Y. M. Chi, A. Ojeda, T. Kerth, S. Makeig,
T.-P. Jung, and G. Cauwenberghs, ``Real-time neuroimaging and cognitive monitoring
using wearable dry EEG,'' \textit{IEEE Transactions on Biomedical Engineering},
vol. 62, no. 11, pp. 2553--2567, 2015.

\bibitem{kappel2019earEEG}
S. L. Kappel, M. L. Rank, H. O. Toft, M. Andersen, and P. Kidmose,
``Dry-contact electrode ear-EEG,'' \textit{IEEE Transactions on Biomedical
Engineering}, vol. 66, no. 1, pp. 150--158, Jan. 2019.

\bibitem{liu2024mi}
H. Liu, P. Wei, H. Wang, X. Lv, W. Duan, M. Li, Y. Zhao, Q. Wang, X. Chen, G. Shi,
B. Han, and J. Hao, ``An EEG motor imagery dataset for brain computer interface in
acute stroke patients,'' \textit{Scientific Data}, vol. 11, no. 1, p. 131, 2024,
doi: 10.1038/s41597-023-02787-8.

\bibitem{leeb2007virtual}
R. Leeb, F. Lee, C. Keinrath, R. Scherer, H. Bischof, and G. Pfurtscheller,
``Brain-computer communication: motivation, aim, and impact of exploring a virtual
apartment,'' \textit{IEEE Transactions on Neural Systems and Rehabilitation Engineering},
vol. 15, no. 4, pp. 473--482, 2007.

\bibitem{tangermann2012review}
M. Tangermann, K.-R. M\"{u}ller, A. Aertsen, N. Birbaumer, C. Braun, C. Brunner,
R. Leeb, C. Mehring, K. J. Miller, G. R. M\"{u}ller-Putz, G. Nolte, G. Pfurtscheller,
H. Preissl, G. Schalk, A. Schl\"{o}gl, C. Vidaurre, S. Waldert, and B. Blankertz,
``Review of the BCI competition IV,'' \textit{Frontiers in Neuroscience}, vol. 6,
art. 55, 2012, doi: 10.3389/fnins.2012.00055.

\bibitem{ang2012fbcsp}
K. K. Ang, Z. Y. Chin, C. Wang, C. Guan, and H. Zhang, ``Filter bank common spatial
pattern algorithm on BCI competition IV datasets 2a and 2b,'' \textit{Frontiers in
Neuroscience}, vol. 6, art. 39, 2012, doi: 10.3389/fnins.2012.00039.

\bibitem{pfurtscheller1999erd}
G. Pfurtscheller and F. H. Lopes da Silva, ``Event-related EEG/MEG synchronization
and desynchronization: basic principles,'' \textit{Clinical Neurophysiology},
vol. 110, no. 11, pp. 1842--1857, 1999.

\bibitem{ramoser2000csp}
H. Ramoser, J. M\"{u}ller-Gerking, and G. Pfurtscheller, ``Optimal spatial filtering
of single trial EEG during imagined hand movement,'' \textit{IEEE Transactions on
Rehabilitation Engineering}, vol. 8, no. 4, pp. 441--446, 2000.

\bibitem{cortes1995svm}
C. Cortes and V. Vapnik, ``Support-vector networks,'' \textit{Machine Learning},
vol. 20, no. 3, pp. 273--297, 1995.

\bibitem{barachant2012riemann}
A. Barachant, S. Bonnet, M. Congedo, and C. Jutten, ``Multiclass brain-computer
interface classification by Riemannian geometry,'' \textit{IEEE Transactions on Biomedical
Engineering}, vol. 59, no. 4, pp. 920--928, 2012.

\bibitem{xu2023review}
C. Xu and R.-Z. Xia, ``EEG signal classification and feature extraction methods based
on deep learning: A review,'' in \textit{Proc. 2023 2nd International Conference on
Big Data, Information and Computer Network (BDICN)}, Xishuangbanna, China, 2023,
pp. 203--206, doi: 10.1109/BDICN58493.2023.00046.

\bibitem{song2023conformer}
Y. Song, Q. Zheng, B. Liu, and X. Gao, ``EEG Conformer: Convolutional transformer for
EEG decoding and visualization,'' \textit{IEEE Transactions on Neural Systems and
Rehabilitation Engineering}, vol. 31, pp. 710--719, 2023,
doi: 10.1109/TNSRE.2022.3230250.

\bibitem{jerrin2021imagined}
J. T. Panachakel and A. G. Ramakrishnan, ``Decoding imagined speech from EEG using
transfer learning,'' \textit{IEEE Access}, vol. 9, pp. 135371--135383, 2021,
doi: 10.1109/ACCESS.2021.3116196.

\bibitem{su2024hybrid}
J. Su, J. Wang, W. Wang, Y. Wang, C. Bunterngchit, P. Zhang, and Z.-G. Hou,
``An adaptive hybrid brain-computer interface for hand function rehabilitation
of stroke patients,'' \textit{IEEE Transactions on Neural Systems and
Rehabilitation Engineering}, vol. 32, pp. 2950--2960, 2024.

\bibitem{krishna2019nospeech}
G. Krishna, C. Tran, J. Yu, and A. H. Tewfik, ``Speech recognition with no speech
or with noisy speech,'' in \textit{Proc. IEEE International Conference on Acoustics,
Speech and Signal Processing (ICASSP)}, pp. 1090--1094, 2019.

\bibitem{krishna2020eeg}
G. Krishna, C. Tran, Y. Han, M. Carnahan, and A. H. Tewfik, ``Speech synthesis
using EEG,'' in \textit{Proc. IEEE International Conference on Acoustics, Speech
and Signal Processing (ICASSP)}, pp. 1235--1239, 2020.

\bibitem{zhao2015phonological}
S. Zhao and F. Rudzicz, ``Classifying phonological categories in imagined and
articulated speech,'' in \textit{Proc. IEEE International Conference on Acoustics,
Speech and Signal Processing (ICASSP)}, Brisbane, Australia, Apr. 2015,
pp. 992--996.

\bibitem{bollens2025auditory}
L. Bollens, C. Piffay, B. Accou, J. Vanthornhout, H. Van Hamme, and T. Francart,
``Auditory EEG decoding challenge for ICASSP 2024,'' \textit{IEEE Open Journal
of Signal Processing}, vol. 6, pp. 1--13, 2025,
doi: 10.1109/OJSP.2025.3534122.

\bibitem{bird2018mental}
J. J. Bird, A. Ek\'{a}rt, L. J. Manso, D. R. Faria, and E. P. Ribeiro, ``A study
on mental state classification using EEG-based brain-machine interface,'' in
\textit{Proc. 2018 International Conference on Intelligent Systems (IS)},
pp. 1--6, 2018.

\bibitem{jawed2024visual}
S. Jawed, I. Faye, and A. S. Malik, ``Deep learning-based assessment model for
real-time identification of visual learners using raw EEG,'' \textit{IEEE
Transactions on Neural Systems and Rehabilitation Engineering}, vol. 32,
pp. 378--390, 2024, doi: 10.1109/TNSRE.2024.3351694.

\bibitem{ming2021drowsiness}
Y. Ming, D. Wu, Y.-K. Wang, Y. Shi, and C.-T. Lin, ``EEG-based drowsiness
estimation for driving safety using deep Q-learning,'' \textit{IEEE Transactions
on Emerging Topics in Computational Intelligence}, vol. 5, no. 4, pp. 583--594,
2021, doi: 10.1109/TETCI.2020.2997031.

\bibitem{jiri2025drowsiness}
E. Jiri, A. Maraoui, I. Werda, ``Driver drowsiness detection using boosting
algorithms with integrated EEG spectral entropy and PSD features,'' in
\textit{Proc. 2025 IEEE 22nd International Conference on Sciences and Techniques
of Automatic Control and Computer Engineering (STA)}, pp. 253--257, 2025.

%====================== Chapters 3 to 5 : System Overview, Hardware and Software ======================

\bibitem{jasper1958tentwenty}
H. H. Jasper, ``The ten-twenty electrode system of the International Federation,''
\textit{Electroencephalography and Clinical Neurophysiology}, vol. 10,
pp. 371--375, 1958.

% TODO: replace the URL and access date below with the exact page you used.
\bibitem{gelo2026manual}
GELO, ``GELO BCI Digital EEG Recorder Headphones: user manual and SDK
documentation,'' GELO, 2026. [Online]. Available: \url{https://www.gelo.one/}.
[Accessed: April 10, 2026].

\bibitem{rpi5brief}
Raspberry Pi Ltd., ``Raspberry Pi 5 product brief,'' 2023. [Online]. Available:
\url{https://datasheets.raspberrypi.com/rpi5/raspberry-pi-5-product-brief.pdf}.
[Accessed: April 10, 2026].

% TODO: confirm the exact model of your panel and update the wiki link if needed.
\bibitem{waveshare7inch}
Waveshare, ``7inch HDMI LCD capacitive touch display,'' Waveshare Wiki. [Online].
Available: \url{https://www.waveshare.com/wiki/7inch_HDMI_LCD}.
[Accessed: April 10, 2026].

\bibitem{gramfort2013mne}
A. Gramfort, M. Luessi, E. Larson, D. A. Engemann, D. Strohmeier, C. Brodbeck,
R. Goj, M. Jas, T. Brooks, L. Parkkonen, and M. S. H\"{a}m\"{a}l\"{a}inen,
``MEG and EEG data analysis with MNE-Python,'' \textit{Frontiers in Neuroscience},
vol. 7, art. 267, 2013.

\bibitem{pedregosa2011sklearn}
F. Pedregosa, G. Varoquaux, A. Gramfort, V. Michel, B. Thirion, O. Grisel,
M. Blondel, P. Prettenhofer, R. Weiss, V. Dubourg, J. Vanderplas, A. Passos,
D. Cournapeau, M. Brucher, M. Perrot, and \'{E}. Duchesnay, ``Scikit-learn:
Machine learning in Python,'' \textit{Journal of Machine Learning Research},
vol. 12, pp. 2825--2830, 2011.

%====================== Chapter 6 : Results ======================

\bibitem{welch1967}
P. Welch, ``The use of fast Fourier transform for the estimation of power spectra:
A method based on time averaging over short, modified periodograms,''
\textit{IEEE Transactions on Audio and Electroacoustics}, vol. 15, no. 2,
pp. 70--73, 1967.

\bibitem{he2020ea}
H. He and D. Wu, ``Transfer learning for brain-computer interfaces: A Euclidean space
data alignment approach,'' \textit{IEEE Transactions on Biomedical Engineering}, vol. 67,
no. 2, pp. 399--410, 2020.

\end{thebibliography}
